\documentclass[12pt,a4paper]{report}
% Общая преамбула. Подключается из subject/main.tex после \documentclass.

% ============================================================
% Базовые шрифты и кодировка
% ============================================================
\usepackage[T2A]{fontenc}
\usepackage[utf8]{inputenc}
\usepackage[russian]{babel}

\renewcommand{\rmdefault}{cmr}
\renewcommand{\sfdefault}{cmss}
\renewcommand{\ttdefault}{cmtt}
\renewcommand{\familydefault}{\rmdefault}

% ============================================================
% Математика и базовое оформление
% ============================================================
\usepackage{amsmath,amssymb,amsthm,mathtools}
\usepackage{mathrsfs}
\usepackage{bbold}
\usepackage{enumitem}
\usepackage{xparse}
\usepackage{xcolor}
\usepackage{microtype}
\usepackage{geometry}          % для настройки полей
\usepackage{parskip}           % для расстояний между абзацами (без параметров)
\usepackage{hyperref}
\usepackage[nameinlink,noabbrev]{cleveref}
\usepackage{tikz}
\usetikzlibrary{
    positioning,
    arrows.meta,
    calc,
    intersections,
    patterns,
    shapes.geometric,
    decorations.markings,
    matrix,
    fit,
    backgrounds
}

% ============================================================
% Рисунки, графики, схемы, таблицы и фрагменты кода
% ============================================================
\usepackage{graphicx}      % изображения
\usepackage{float}         % точное размещение рисунков/таблиц через [H]
\usepackage{caption}       % подписи к рисункам и таблицам
\usepackage{subcaption}    % несколько рисунков внутри одной фигуры
\usepackage{wrapfig}       % обтекание рисунков текстом
\usepackage{pgfplots}      % графики функций и численных данных
\pgfplotsset{compat=1.18}
\usepackage{tikz-cd}       % коммутативные диаграммы

\usepackage{booktabs}      % аккуратные таблицы
\usepackage{array}
\usepackage{tabularx}
\usepackage{longtable}
\usepackage{multirow}

\usepackage{listings}      % код: C/C++/Python/SQL и др.

% Дополнительная математика
\usepackage{bm}            % жирные математические символы: \bm{x}
\usepackage{bbm}           % красивый blackboard bold для цифр: \mathbbm{0}, \mathbbm{1}
\usepackage{esint}         % \oiint, \varoiint и другие интегралы
\usepackage{cancel}        % \cancel{}, \bcancel{}, \xcancel{}

% ============================================================
% Отступы и расстояния
% ============================================================
\geometry{a4paper,top=2cm,bottom=2cm,left=3cm,right=3cm,marginparwidth=1.75cm}
% parskip уже загружен без параметров, что даёт нулевой отступ первой строки
% и положительный межстрочный интервал между абзацами.

\setlength{\emergencystretch}{3em}

\hypersetup{
    colorlinks=true,
    linkcolor=black,
    urlcolor=black,
    citecolor=black,
    linktoc=all
}

% ============================================================
% Доказательства
% ============================================================
\addto\captionsrussian{%
    \renewcommand{\proofname}{Доказательство}%
}

\makeatletter
\renewenvironment{proof}[1][\proofname]{%
    \par
    \begingroup
    \leftskip=\dimexpr-\@totalleftmargin\relax
    \rightskip=0pt
    \parindent=0pt
    \noindent\textit{\underline{#1.}}\par\nobreak
    \noindent\ignorespaces
}{%
    \unskip\nobreak\hspace{0.15em}\ensuremath{\blacksquare}%
    \par
    \endgroup
}
\makeatother

% ============================================================
% Заголовки
% ============================================================
\usepackage{titlesec}
\titleclass{\chapter}{straight}

\titleformat{\chapter}[display]
  {\normalfont\huge\bfseries\raggedright}
  {\chaptername\ \thechapter}
  {0.7em}
  {\Huge\raggedright}

\titleformat{\section}
  {\normalfont\Large\bfseries}
  {\thesection}{0.7em}{}

\titleformat{\subsection}
  {\normalfont\large\bfseries}
  {\thesubsection}{0.7em}{}

\titleformat{\subsubsection}
  {\normalfont\normalsize\bfseries}
  {\thesubsubsection}{0.7em}{}

\titlespacing*{\chapter}{0pt}{0.45em}{0.45em}
\titlespacing*{\section}{0pt}{0.55em}{0.25em}
\titlespacing*{\subsection}{0pt}{0.45em}{0.18em}
\titlespacing*{\subsubsection}{0pt}{0.35em}{0.15em}

% ============================================================
% Лекции
% ============================================================
\newcounter{lection}
\NewDocumentCommand{\newlection}{m}{%
    \refstepcounter{lection}%
    \par\smallskip
    {\centering\color{gray!70}\normalsize Лекция \thelection. #1\par}
    \vspace{0.1em}
}

% ============================================================
% Теоремы, определения, замечания
% ============================================================
\makeatletter
\NewDocumentCommand{\definition}{o +m}{%
    \par
    \noindent
    \ifdim\@totalleftmargin>0pt
        \hspace*{\dimexpr-\@totalleftmargin\relax}%
    \fi
    \textbf{Определение\IfValueT{#1}{ (#1)}.}~#2
    \par
}
\makeatother

\ExplSyntaxOn

\NewDocumentCommand{\theorem}{o +m}{%
    \par\noindent
    \textbf{%
        Теорема
        \IfValueT{#1}{%
            \regex_match:nnTF { \A \d+ \Z } { #1 }
                {~#1}
                {~(#1)}
        }%
        .
    }~#2
    \par
}

\ExplSyntaxOff

\NewDocumentCommand{\proposal}{o +m}{%
    \par\noindent
    \textbf{Предложение\IfValueT{#1}{~#1}.}~#2
    \par
}

\NewDocumentCommand{\fact}{o +m}{%
    \par\noindent
    \textit{Факт\IfValueT{#1}{ (#1)}.}~#2
    \par
}

\NewDocumentCommand{\question}{+m}{%
    \par\noindent
    \textbf{Вопрос.}~#1
    \par
}

\NewDocumentCommand{\note}{+m}{%
    \par\noindent
    \textbf{Замечание.}~#1
    \par
}

\NewDocumentCommand{\corollary}{o +m}{%
    \par\noindent
    \textbf{Следствие\IfValueT{#1}{~#1}.}~#2
    \par
}

\NewDocumentCommand{\example}{+m}{%
    \par\noindent
    \textbf{Пример.}~#1
    \par
}

% ============================================================
% Списки
% ============================================================
\NewDocumentCommand{\numbers}{+m}{%
    \begin{enumerate}[leftmargin=2em,itemsep=0.25em,topsep=0.3em]
    #1
    \end{enumerate}
}

\NewDocumentCommand{\bullets}{+m}{%
    \begin{itemize}[leftmargin=2em,itemsep=0.25em,topsep=0.3em]
    #1
    \end{itemize}
}

\NewDocumentCommand{\provehere}{+m}{%
    \par\smallskip
    \begin{proof}#1\end{proof}
}

\NewDocumentCommand{\provewthen}{m +m}{%
    \par\smallskip
    \begin{proof}[#1]#2\end{proof}
}

\NewDocumentCommand{\provebullets}{+m}{%
    \begin{proof}
    \begin{itemize}[leftmargin=2em,itemsep=0.25em,topsep=0.3em]
    #1
    \end{itemize}
    \end{proof}
}

\NewDocumentCommand{\provenumbers}{+m}{%
    \begin{proof}
    \begin{enumerate}[leftmargin=2em,itemsep=0.25em,topsep=0.3em]
    #1
    \end{enumerate}
    \end{proof}
}

% ============================================================
% Дополнительные команды
% ============================================================
\usepackage{empheq}      % для \encircle
\usepackage{changepage}  % для \blockindent

% Леммы
\newtheorem{lemma_env}{Лемма}[section]
\newcommand{\lemma}[2][]{\begin{lemma_env}[#1]#2\end{lemma_env}}

% Отступ (переименовано из \indent, чтобы не ломать ядро LaTeX)
\newcommand{\blockindent}[1]{%
  \begin{adjustwidth}{1.5cm}{}#1\end{adjustwidth}%
}

% Одностраничные блоки
\newcommand{\singlepage}[1]{\begin{samepage}#1\end{samepage}}

% Сборки формул
\renewcommand{\gather}[1]{\begin{gather*}#1\end{gather*}}
\renewcommand{\multline}[1]{\begin{multline*}#1\end{multline*}}

% Обрамлённые формулы
\newcommand{\encircle}[1]{\begin{empheq}[box=\fbox]{align*}#1\end{empheq}}

% Математические операторы и символы
\newcommand{\proddots}{\cdot\ldots\cdot}
\newcommand{\abs}[1]{\left|#1\right|}
\newcommand{\angles}[1]{\left\langle#1\right\rangle}
\newcommand{\der}[2]{\frac{\partial #1}{\partial #2}}
\DeclareMathOperator{\differential}{d}
\newcommand{\dif}{\differential\!}
\newcommand{\fmax}{\lor}
\newcommand{\fmin}{\land}
\DeclareMathOperator*{\esssup}{ess\,sup}
\DeclareMathOperator*{\essssup}{ess\,sup}
\DeclareMathOperator*{\esssssup}{ess\,sup}
\newcommand{\Map}{\longrightarrow}
\newcommand{\almost}[1]{\underset{#1}{\overset{\textup{почти всюду}}\Map}}
\newcommand{\circlesign}[1]{\mathrel{\tikz[baseline=-0.5ex]{\node[draw, circle, inner sep=1pt] (char) {\vphantom{$-$}$#1$};}}}

% Дополнительные команды для доказательств
\newcommand{\indentlemma}[3][]{\blockindent{\lemma[#1]{#2\prove[Доказательство леммы]{#3}}}}
\newcommand{\prove}[2][Доказательство]{\begin{proof}[#1]#2\end{proof}}

% ============================================================
% Часто используемые математические обозначения
% ============================================================
\newcommand{\N}{\mathbb{N}}
\newcommand{\Z}{\mathbb{Z}}
\newcommand{\Q}{\mathbb{Q}}
\newcommand{\R}{\mathbb{R}}
\renewcommand{\C}{\mathbb{C}}
\newcommand{\Ord}{\mathrm{Ord}}
\newcommand{\Card}{\mathrm{Card}}

\renewcommand{\o}{\varnothing}
\newcommand{\bs}{\setminus}
\newcommand{\inv}{^{\complement}}
\newcommand{\then}{\Rightarrow}
\newcommand{\map}{\to}
\newcommand{\bydef}{\coloneqq}
\newcommand{\all}[1]{\left\{\begin{aligned}#1\end{aligned}\right.}
\newcommand{\any}[1]{\left[\begin{aligned}#1\end{aligned}\right.}
\newcommand{\exist}{\exists}
\newcommand{\dom}{\operatorname{dom}}
\newcommand{\rng}{\operatorname{rng}}
\newcommand{\defset}[2]{\left\{#1\;\middle|\;#2\right\}}
\newcommand{\switch}[1]{\left\{\begin{aligned}#1\end{aligned}\right.}

% ------------------------------------------------------------
% Большие операторы: индексы автоматически ставятся сверху/снизу
% ------------------------------------------------------------
\let\oldlim\lim
\renewcommand{\lim}{\oldlim\limits}
\let\oldsup\sup
\renewcommand{\sup}{\oldsup\limits}
\let\oldinf\inf
\renewcommand{\inf}{\oldinf\limits}
\let\oldmax\max
\renewcommand{\max}{\oldmax\limits}
\let\oldmin\min
\renewcommand{\min}{\oldmin\limits}

\let\oldsum\sum
\renewcommand{\sum}{\oldsum\limits}
\let\oldprod\prod
\renewcommand{\prod}{\oldprod\limits}
\let\oldcoprod\coprod
\renewcommand{\coprod}{\oldcoprod\limits}

\let\oldbigcap\bigcap
\renewcommand{\bigcap}{\oldbigcap\limits}
\let\oldbigcup\bigcup
\renewcommand{\bigcup}{\oldbigcup\limits}
\let\oldbigsqcup\bigsqcup
\renewcommand{\bigsqcup}{\oldbigsqcup\limits}

\let\oldbigvee\bigvee
\renewcommand{\bigvee}{\oldbigvee\limits}
\let\oldbigwedge\bigwedge
\renewcommand{\bigwedge}{\oldbigwedge\limits}
\let\oldbigoplus\bigoplus
\renewcommand{\bigoplus}{\oldbigoplus\limits}
\let\oldbigotimes\bigotimes
\renewcommand{\bigotimes}{\oldbigotimes\limits}
\let\oldbigodot\bigodot
\renewcommand{\bigodot}{\oldbigodot\limits}
\let\oldbiguplus\biguplus
\renewcommand{\biguplus}{\oldbiguplus\limits}

% Красивые blackboard-bold цифры (\mathbb не во всех шрифтах содержит цифры)
\newcommand{\bbzero}{\mathbbm{0}}
\newcommand{\bbone}{\mathbbm{1}}

% Универсальные обозначения для разных курсов
\newcommand{\norm}[1]{\left\lVert #1 \right\rVert}
\newcommand{\floor}[1]{\left\lfloor #1 \right\rfloor}
\newcommand{\ceil}[1]{\left\lceil #1 \right\rceil}
\newcommand{\paren}[1]{\left( #1 \right)}
\newcommand{\brackets}[1]{\left[ #1 \right]}
\newcommand{\set}[1]{\left\{ #1 \right\}}
\newcommand{\conj}[1]{\overline{#1}}
\newcommand{\vect}[1]{\bm{#1}}

\DeclareMathOperator{\RePart}{Re}
\DeclareMathOperator{\ImPart}{Im}
\DeclareMathOperator{\Arg}{Arg}
\DeclareMathOperator{\Ker}{Ker}
\DeclareMathOperator{\Image}{Im}
\DeclareMathOperator{\rank}{rank}
\DeclareMathOperator{\tr}{tr}
\DeclareMathOperator{\sgn}{sgn}
\DeclareMathOperator{\lcm}{lcm}
\DeclareMathOperator{\Span}{span}

% Differential notation in math; preserve the text dot-below accent.
\let\textdotbelow\d
\DeclareRobustCommand{\d}{\ifmmode\mathrm{d}\else\expandafter\textdotbelow\fi}

\endinput

% Всё ниже оставлено из исходного файла как справочный образец титульной
% страницы и LaTeX не читает после команды \endinput.
% ============================================================
% Информация о документе и титульная страница
% ============================================================
\date{III семестр, 2026 г.}
\title{Название предмета. Неофициальный конспект}
\author{Лектор: ФИО лектора\\
Конспект: Викторова Дарья, 25.Б03-пу}

\begin{document}

\begin{titlepage}
    \thispagestyle{empty}
    \begin{center}
        \vspace*{0.6cm}

        {\large Санкт-Петербургский государственный университет\par}
        \vspace{0.35em}
        {\normalsize Факультет прикладной математики~--- процессов управления\par}

        \vspace*{5.1cm}

        {\fontsize{30}{35}\selectfont\bfseries Название предмета\par}
        \vspace{0.65em}
        {\fontsize{17}{21}\selectfont\itshape Неофициальный конспект\par}

        \vspace{2.7cm}

        {\normalsize
        \begin{tabular}{@{}rl@{}}
            \textit{Лектор:} & ФИО лектора \\
            \textit{Конспект:} & Викторова Дарья, 25.Б03-пу
        \end{tabular}\par}

        \vfill

        {\large III семестр\par}
        \vspace{0.22em}
        {\large 2026 год\par}
        \vspace*{0.55cm}
    \end{center}
\end{titlepage}

\tableofcontents
\newpage

\setcounter{lection}{0}
\newlection{Число Месяц 2026 г.}

\end{document}



\title{Математический анализ. Неофициальный конспект}
\author{Лектор: Скопина Мария Александровна\\
Конспект: Викторова Дарья, 25.Б03-пу}
\date{III семестр, 2026 г.}

\begin{document}

\begin{titlepage}
    \thispagestyle{empty}
    \begin{center}
        \vspace*{0.6cm}

        {\large Санкт-Петербургский государственный университет\par}
        \vspace{0.35em}
        {\normalsize Факультет прикладной математики~--- процессов управления\par}

        \vspace*{5.1cm}

        {\fontsize{30}{35}\selectfont\bfseries Математический анализ\par}
        \vspace{0.65em}
        {\fontsize{17}{21}\selectfont\itshape Неофициальный конспект\par}

        \vspace{2.7cm}

        {\normalsize
        \begin{tabular}{@{}rl@{}}
            \textit{Лектор:} & ФИО лектора \\
            \textit{Конспект:} & Викторова Дарья, 25.Б03-пу
        \end{tabular}\par}

        \vfill

        {\large III семестр\par}
        \vspace{0.22em}
        {\large 2026 год\par}
        \vspace*{0.55cm}
    \end{center}
\end{titlepage}

\tableofcontents
\newpage

\setcounter{lection}{0}
\newlection{2 сентября 2026 г.}

\chapter{Дифференциальное исчисление функций многих переменных}

\section{Многомерные пространства}

\definition{
    $\mathbb{R}^d$ --- \textit{евклидово вещественное пространство}, \\
    $x=(x_1,\ldots,x_d)$, $x_k\in\mathbb{R}$, --- элементы евклидова вещественного пространства.
}

Введём операции:
\[
    x\pm y:=(x_1\pm y_1,\ldots,x_d\pm y_d),
    \qquad
    \alpha\in\mathbb{R}=\mathbb{R}^1\colon
    \alpha x=(\alpha x_1,\ldots,\alpha x_d).
\]

Евклидово пространство является \textit{векторным пространством}.

\subsection{Аксиомы векторного пространства}

\begin{enumerate}
    \item Коммутативность сложения: $\forall x,y\in\mathbb{R}^d\colon x+y=y+x$.
    \item Ассоциативность сложения: $\forall x,y,z\in\mathbb{R}^d\colon (x+y)+z=x+(y+z)$.
    \item Существование нейтрального элемента по сложению: $\exists\,\mathbb{0}\in\mathbb{R}^d\colon \forall x\in\mathbb{R}^d\colon x+\mathbb{0}=x$, $\mathbb{0}=(0,\ldots,0)$.
    \item $\forall x\in\mathbb{R}^d\ \exists(-x)\colon x+(-x)=\mathbb{0}$.
    \item Ассоциативность умножения на скаляр: $\alpha\cdot(\beta x)=(\alpha\beta)\cdot x$.
    \item $(\alpha+\beta)\cdot x=\alpha x+\beta x$.
    \item $1\cdot x=x$.
\end{enumerate}

\definition{
    $\rho(x,y):=
    \sqrt{(x_1-y_1)^2+\ldots+(x_d-y_d)^2}
    =\rho(x-y,\mathbb{0})$
    --- \textit{метрика} (\textit{расстояние}).
}

\subsection{Свойства метрики}

\begin{enumerate}
    \item $\rho(x,y)=0\iff x=y$.
    \item $\rho(x,y)\leq\rho(x,z)+\rho(z,y)\quad \forall x,y,z\in\mathbb{R}^d$.

\begin{proof}
    \noindent
        $\displaystyle
        \rho^2(x,y)
        =(x_1-y_1)^2+\ldots+(x_d-y_d)^2=$%

    \noindent
       $=((x_1-z_1)+(z_1-y_1))^2+\ldots+((x_d-z_d)+(z_d-y_d))^2=
        $%

    \noindent
        $\displaystyle
        =\sum_{k=1}^d (x_k-z_k)^2+\sum_{k=1}^d(z_k-y_k)^2
        +2\sum_{k=1}^d(x_k-z_k)(z_k-y_k)\leq
        $%

    \noindent
        $\displaystyle
        \leq \rho^2(x,z)+\rho^2(z,y)+2\rho(x,z)\rho(z,y)
        =(\rho(x,z)+\rho(z,y))^2.
        $%

    Проверим последнее неравенство для $d=2$. Обозначим $x_k-z_k=a_k$, $z_k-y_k=b_k$. \\
    Надо показать, что $\displaystyle a_1b_1+a_2b_2 \leq \sqrt{a_1^2+a_2^2}\cdot \sqrt{b_1^2+b_2^2}.$%

    \noindent
        $\displaystyle
        (a_1b_1+a_2b_2)^2
        \leq(a_1^2+a_2^2)(b_1^2+b_2^2),$

    \noindent
        $\displaystyle
        (a_1b_1)^2+2a_1b_1a_2b_2+(a_2b_2)^2
        \leq a_1^2b_1^2+a_1^2b_2^2+a_2^2b_1^2+a_2^2b_2^2,$

    \noindent
        $\displaystyle
        2a_1b_2a_2b_1
        \leq a_1^2b_2^2+a_2^2b_1^2.$
        
\end{proof}
\end{enumerate}

\definition{
    Рассмотрим $x,y\in\mathbb{R}^d$. $\langle x,y\rangle:=x_1y_1+\ldots+x_dy_d$ --- \textit{скалярное произведение}.
}

\subsection{Свойства скалярного произведения}

\begin{enumerate}
    \item $\forall x,y\in\mathbb{R}^d\colon \langle x,y\rangle=\langle y,x\rangle$.
    \begin{proof}
        По определению.
    \end{proof}

    \item $\forall x,x',y\in\mathbb{R}^d\colon \langle x+x',y\rangle=\langle x,y\rangle+\langle x',y\rangle$.
    \begin{proof}
        По определению.
    \end{proof}

    \item $\forall \lambda\in\mathbb{R},\ \forall x,y\in\mathbb{R}^d\colon \langle\lambda x,y\rangle=\langle x,\lambda y\rangle$.
    \begin{proof}
        По определению.
    \end{proof}

    \item $\forall x\in\mathbb{R}^d\colon \langle x,x\rangle\geq 0$.

    \definition{
        $\lVert x\rVert:=\sqrt{\langle x,x\rangle}$ --- \textit{норма}.
    }

    \item $\langle x,x\rangle=0\iff x=\mathbb{0}$.

    \item $\forall x,y\in\mathbb{R}^d\colon |\langle x,y\rangle|\leq\lVert x\rVert\cdot\lVert y\rVert$.
    \begin{proof}
        Возьмём $\lambda\in\mathbb{R}\colon \langle\lambda x+y,\lambda x+y\rangle\geq 0$. Возьмём
        $\displaystyle \lambda=-\frac{\langle x,y\rangle}{\lVert x\rVert^2}$. Тогда
        \[
        \displaystyle \lambda^2\langle x,x\rangle +2\lambda\langle x,y\rangle +\langle y,y\rangle\geq 0.
        \]

        Следовательно,

        \noindent
        $\displaystyle
            \frac{(\langle x,y\rangle)^2}{\lVert x\rVert^2}
            -2\frac{(\langle x,y\rangle)^2}{\lVert x\rVert^2}
            +\langle y,y\rangle\geq 0,$

        \noindent
            $\displaystyle
            (\langle x,y\rangle)^2
            -\lVert x\rVert^2\lVert y\rVert^2\leq 0,
            $

        \noindent
            $\displaystyle
            (\langle x,y\rangle)^2
            \leq
            \lVert x\rVert^2\lVert y\rVert^2.
            $
    \end{proof}
\end{enumerate}

\proposal[1]{
    $\forall x,y\in\mathbb{R}^d\colon
    \lVert x+y\rVert\leq\lVert x\rVert+\lVert y\rVert$.
}

\begin{proof}
    \[
        \lVert x+y\rVert
        =\sqrt{\langle x+y,x+y\rangle}
        =\sqrt{\langle x,x\rangle+2\langle x,y\rangle+\langle y,y\rangle}
        \leq\sqrt{\lVert x\rVert^2+2\lVert x\rVert\lVert y\rVert+\lVert y\rVert^2}=
    \]
    $=\lVert x\rVert+\lVert y\rVert.$
\end{proof}

\corollary{
    $\bigl|\lVert x\rVert-\lVert y\rVert\bigr|\leq\lVert x-y\rVert$.
}

\begin{proof}
    а) Пусть $\lVert x\rVert>\lVert y\rVert$. По Предложению 1:
        \[ \displaystyle
        \lVert x\rVert
        =\lVert x-y+y\rVert
        \leq\lVert x-y\rVert+\lVert y\rVert \Rightarrow 
    \lVert x\rVert-\lVert y\rVert
    \leq\lVert x-y\rVert.
    \]

    б) Пусть $\lVert x\rVert<\lVert y\rVert$. Аналогично:
    \[
        \displaystyle
        \lVert y\rVert
        =\lVert y-x+x\rVert
        \leq\lVert y-x\rVert+\lVert x\rVert
        =\lVert x-y\rVert+\lVert x\rVert \Rightarrow \displaystyle
    \lVert y\rVert-\lVert x\rVert
    \leq\lVert x-y\rVert.
    \]

    Следовательно,
    $\displaystyle
    \bigl|\lVert x\rVert-\lVert y\rVert\bigr|
    \leq\lVert x-y\rVert.
    $
\end{proof}

Пусть $x^0\in\mathbb{R}^d$, $\varepsilon>0$. Введём

\numbers{
    \item $U_\varepsilon(x^0):=
    \{x\in\mathbb{R}^d\colon\rho(x,x^0)<\varepsilon\}$
    --- \textit{$\varepsilon$-окрестность точки $x^0$};

    \item $\mathring{U}_\varepsilon(x^0):=
    U_\varepsilon(x^0)\setminus\{x^0\}$
    --- \textit{проколотая $\varepsilon$-окрестность точки $x^0$}.
}

\definition{
    Пусть $E\subset\mathbb{R}^d$, $x^0\in\mathbb{R}^d$.
    Точка $x^0$ называется \textit{предельной точкой} множества $E$, если
    $\forall\varepsilon>0\colon
    E\cap\mathring{U}_\varepsilon(x^0)\neq\varnothing$.
}

Рассмотрим $\{x^{(n)}\}_{n=1}^{\infty}$ --- последовательность в $\mathbb{R}^d$.

\definition{
    Последовательность $\{x^{(n)}\}_{n=1}^{\infty}$ называется
    \textit{сходящейся к точке} $x^0\in\mathbb{R}^d$, если
    $\rho(x^0,x^{(n)})\xrightarrow[n\to\infty]{}0$.

    \textit{Обозначение}:
    $x^{(n)}\to x^0$ или
    $\displaystyle\lim_{n\to\infty}x^{(n)}=x^0$.
}

\theorem[1]{
    Для того чтобы $x^{(n)}\to x^0$, необходимо и достаточно, чтобы
    $x_k^{(n)}\to x_k^0$ для всех $k=\overline{1,d}$.
}
\begin{proof}
    \begin{enumerate}
        \item[$\Longrightarrow$] \textit{Необходимость}. 
        \[
        |x_k^{(n)}-x_k^0|\leq \rho(x^{(n)},x^0)=\sqrt{(x_1^{(n)}-x_1^0)^2+\dots+(x_d^{(n)}-x_d^0)^2}.
        \]
        Так как $\rho(x^{(n)},x^0)\to 0$, то $|x_k^{(n)}-x_k^0|\to 0$ для каждого $k$, значит, $x_k^{(n)}\to x_k^0$.

        \item[$\Longleftarrow$] \textit{Достаточность}. Надо доказать, что
        \[
        \rho(x^{(n)},x^0)\leq |x_1^{(n)}-x_1^0|+\dots +|x_d^{(n)}-x_d^0|.
        \]
        Воспользуемся методом математической индукции.
        \begin{enumerate}
            \item \textit{База индукции}. $d=2$.
            \[
            \sqrt{(x_1^{(n)}-x_1^0)^2+(x_2^{(n)}-x_2^0)^2} \leq |x_1^{(n)}-x_1^0|+|x_2^{(n)}-x_2^0|.
            \]
            Возведём обе части неравенства в квадрат:
            \[
            (x_1^{(n)}-x_1^0)^2+(x_2^{(n)}-x_2^0)^2
            \leq (x_1^{(n)}-x_1^0)^2+(x_2^{(n)}-x_2^0)^2+2|x_1^{(n)}-x_1^0|\cdot|x_2^{(n)}-x_2^0|.
            \]
            База доказана.

            \item $d-1\to d$. По предположению индукции:
            \[
            \sqrt{(x_2^{(n)}-x_2^0)^2+\dots+(x_d^{(n)}-x_d^0)^2}
            \leq |x_2^{(n)}-x_2^0|+\dots+|x_d^{(n)}-x_d^0|.
            \]
            Тогда
            \[
            \rho(x^{(n)},x^0)\leq \sqrt{(x_1^{(n)}-x_1^0)^2+\bigl(|x_2^{(n)}-x_2^0|+\dots + |x_d^{(n)}-x_d^0|\bigr)^2}.
            \]
            Применив неравенство $\sqrt{a^2+b^2}\le a+b$, $a,b\ge0$, получим
            \[
            \rho(x^{(n)},x^0)\le |x_1^{(n)}-x_1^0|+|x_2^{(n)}-x_2^0|+\dots+|x_d^{(n)}-x_d^0|.
            \]
        \end{enumerate}
    \end{enumerate}
\end{proof}

\newlection{04 сентября 2026 г.}

\definition{Множество $E\subset \mathbb R^d$ называется \textit{ограниченным}, если 
\[ \exists M > 0\colon \forall x\in E\colon \|x\|\leq M \quad (\text{или } \rho(x,\mathbb 0)\leq M). \]
}

\theorem[2]{Для любой ограниченной последовательности $\{x^{(n)}\}_{n=1}^\infty$ существует сходящаяся подпоследовательность $\{x^{(n_k)}\}_{k=1}^\infty$; \textit{обозначим} $\lim_{k\to\infty}x^{(n_k)}=a$.}

\begin{proof}
    $\|x^{(n)}\|\leq M$. \\
    Рассмотрим $\{x_1^{(n)}\}$ --- ограниченная последовательность. Тогда $\exists \{x_1^{(n_k)}\}_{k=1}^\infty$ --- сходится. \\
    Рассмотрим $\{x_2^{(n_k)}\}_{k=1}^\infty$ --- ограниченная последовательность. Тогда $\exists \{x_2^{(n_{k_l})}\}_{l=1}^\infty$ --- сходится. \\
    $\dots$ \\
    Продолжая процесс, получим, что для всех координат последовательность сходится. По теореме 1 подпоследовательности сходятся.
\end{proof}

\definition{Пусть $E\subset \mathbb R^d$, $f$ определена на $E$, $x^0$ --- предельная точка множества $E$. \\
Говорят, что $a\in \mathbb R^d$ является \textit{пределом функции} $f$, и пишут $a=\lim_{x\to x^0}f(x)$, если $\forall\{x^{(n)}\}_{n=1}^\infty \subset E\colon \lim x^{(n)}=x^0, \forall n\colon x^{(n)}\neq x^0$, выполняется $\lim f(x^{(n)})=a$.

\begin{enumerate}
    \item $\lim_{x\to\infty} f(x)=\infty$. Это означает, что для любой последовательности $\{x^{(n)}\}$ такой, что $\lim_{n\to\infty}\|x^{(n)}\|=\infty$, выполняется $\lim_{n\to\infty} f(x^{(n)})=\infty$, то есть
    \[
    \forall M>0\ \exists N\colon \forall n\ge N\colon \rho(f(x^{(n)}),\mathbb 0) > M \quad (\text{или } \|f(x^{(n)})\| \ge M).
    \]

    \item $\lim_{x\to a,\ x\to\infty,\ \pm\infty} f(x)=\infty,\ \pm\infty$.
\end{enumerate}}

\definition{
Пусть $E\subset \mathbb R^d$, $x^0\in E$, $f$ определена на $E$. \\
Говорят, что $f$ \textit{непрерывна в точке} $x^0$, если для любой последовательности $\{x^{(n)}\}_{n=1}^\infty \subset E$ такой, что $\lim_{n\to\infty}x^{(n)}=x^0$, выполняется
\[
\lim_{n\to\infty} f(x^{(n)}) = f(x^0).
\]
}

\example{
    $d = 2$.
    \[
    f(x,y)=\frac{x^2y}{x^4+y^2}, \quad E=\mathbb R^2\setminus\{\mathbb 0\}.
    \]
    Исследуем предел функции при $(x,y)\to(0,0)$.

    1) Возьмём последовательность точек $(h,0)$ при $h\to 0$. Тогда $y=0$, следовательно, $f(h,0)=0$ для всех $h\ne 0$. Значит, предел вдоль этой последовательности равен $0$.

    2) Возьмём последовательность точек $(\sqrt h, h)$ при $h\to 0$. Тогда
    \[
    f(\sqrt h,h)=\frac{(\sqrt h)^2 h}{(\sqrt h)^4 + h^2} = \frac{h^2}{h^2 + h^2} = \frac{1}{2}.
    \]
    Следовательно, предел вдоль этой последовательности равен $\frac{1}{2}$.

    Пределы по разным направлениям получились различными. Значит, предела функции при $(x,y)\to(0,0)$ не существует, и функцию нельзя доопределить в точке $(0,0)$ так, чтобы она была непрерывной.
}

\subsection{Свойства непрерывных функций.}
\begin{enumerate}
    \item $f, g$ определены на $E$, непрерывны в $x^0$. Тогда $\alpha f\pm \beta g$ непрерывна в $x^0$.
    \begin{proof}
        Возьмем $\{x^{(n)}\}\subset E$, $\lim_{n\to\infty}x^{(n)}=x^0$. Обозначим $h=\alpha f\pm \beta g$. \\
        Так как $f$ и $g$ непрерывны, то
        \[
        \lim_{n\to\infty} f(x^{(n)})=f(x^0), \qquad
        \lim_{n\to\infty} g(x^{(n)})=g(x^0).
        \]
        Следовательно,
        \[
        \lim_{n\to\infty} h(x^{(n)})
        =\lim_{n\to\infty}\bigl(\alpha f(x^{(n)})\pm \beta g(x^{(n)})\bigr)
        =\alpha f(x^0)\pm \beta g(x^0)=h(x^0).
        \]
        Значит, $h$ непрерывна в $x^0$.
    \end{proof}

    \item $f, g$ определены на $E$, непрерывны в $x^0$. Тогда $f\cdot g$ непрерывна в $x^0$.

    \item $f, g$ определены на $E$, непрерывны в $x^0$, $x^0\in E$, $g(x^0)\neq 0$. Тогда $\frac{f}{g}$ непрерывна в $x^0$.

    \item $f$ непрерывна в $x^0$, $f(x^0)>0$. Тогда $\exists \delta > 0$ такое, что $f(x) > 0$ для всех $x \in U_\delta(x^0)\cap E$.

        \item Пусть $f_1, \dots, f_N$ определены и непрерывны в $x^0$; $g$ определена на $G\subset \mathbb R^N$; $y^0=(f_1(x^0),\dots,f_N(x^0))\in G$; $g$ непрерывна в $y^0$, имеет смысл суперпозиция $g(f_1,\dots ,f_N)$. Тогда $g(f_1,\dots, f_N)$ непрерывна в $x^0$.
    \begin{proof}
        Возьмем последовательность $\{x^{(n)}\}\subset E,\ \lim_{n\to\infty}x^{(n)}=x^0$. Тогда $\lim_{n\to\infty} f_k(x^{(n)})=f_k(x^0)$ для всех $k=\overline{1,N}$. \\
        Обозначим $y^{(n)}=(f_1(x^{(n)}),\dots,f_N(x^{(n)}))$. Тогда $y^{(n)}\to y^0$. Так как $g$ непрерывна в $y^0$, то
        \[
        \lim_{n\to\infty} g(y^{(n)})=g(y^0)=g(f_1(x^0),\dots,f_N(x^0)).
        \]
        Следовательно, $\lim_{n\to\infty} g(f_1(x^{(n)}),\dots,f_N(x^{(n)}))=g(f_1(x^0),\dots,f_N(x^0))$, то есть суперпозиция непрерывна в $x^0$.
    \end{proof}
\end{enumerate}

\theorem[Вейерштрасса]{Пусть $K$ --- компакт, $f$ непрерывна на $K$. \\
Тогда $f$ ограничена на $K$ и достигает своих наибольшего и наименьшего значений.
}
\begin{proof}
    \begin{enumerate}
        \item Докажем ограниченность. Предположим противное: $f$ не ограничена. Тогда $m=\sup\limits_{x\in K}|f(x)|$.

        \begin{enumerate}
            \item $m = \infty$.
            \[
            \forall M>0\ \exists x=x(M)\colon |f(x)|\geq M.
            \]
            Возьмём точки $x^{(n)}\colon |f(x^{(n)})|>n$. \\
            Рассмотрим $\{x^{(n)}\}$ --- ограничена. Следовательно, покоординатные последовательности $\{x_l^{(n)}\}$ ограничены. \\
            Рассмотрим $\{x_1^{(n)}\}$ --- ограничена. Следовательно, $\exists$ подпоследовательность $\{x^{(n_k)}\}$ такая, что $x_1^{(n_k)}$ сходится, то есть $\lim\limits_{k\to\infty} x_1^{(n_k)}=x_1^0$. \\
            Рассмотрим $\{x_2^{(n_k)}\}$ --- ограничена. Значит, $\exists$ подпоследовательность $\{x^{(n_{k_l})}\}$ такая, что $x_2^{(n_{k_l})}$ сходится. При этом
            \[
            \begin{cases}
               \lim\limits_{l\to\infty} x_2^{(n_{k_l})}=x_2^0,\\
               \lim\limits_{l\to\infty} x_1^{(n_{k_l})}=x_1^0.
            \end{cases}
            \]
            $\dots$ \\
            Продолжая этот процесс, получим, что для подпоследовательности $\{x^{(n_{k_l})}\}$ все координаты сходятся, то есть $x^{(n_{k_l})}\to x^0\in K$ (в силу компактности). \\
            
            Тогда для этой подпоследовательности $|f(x^{(n_{k_l})})|>n_{k_l}\to\infty$, то есть $\lim\limits_{l\to\infty} |f(x^{(n_{k_l})})| = \infty$. \\
            В силу непрерывности $f$ в точке $x^0$ имеем $\lim\limits_{l\to\infty} f(x^{(n_{k_l})}) = f(x^0)$. \\
            
            Следовательно, $\lim\limits_{l\to\infty} |f(x^{(n_{k_l})})| = |f(x^0)|$,
            что является конечным числом.
        \end{enumerate}

        Получили противоречие. Следовательно, $|m|<\infty$, то есть $f$ ограничена.

        \item Докажем достижение максимума. Пусть $m=\sup\limits_{x\in K} f(x)$. Будем считать, что $m>0$. По определению супремума
        \[
        \forall \varepsilon>0\ \exists x=x(\varepsilon)\in K\colon f(x)>m-\varepsilon.
        \]
        Полагая $\varepsilon=\frac1n$, получим $\forall n\in\mathbb N$ точку $x^{(n)}\in K \colon f(x^{(n)})>m-\frac1n$. \\
        
        Возьмем подпоследовательность $\{x^{(n_k)}\}$ --- сходится, $\lim\limits_{k\to\infty} x^{(n_k)}=x^0$. Тогда в силу непрерывности $f$ имеем $\lim\limits_{k\to\infty} f(x^{(n_k)})=f(x^0)$. \\
        С другой стороны, из неравенства $m-\frac{1}{n_k}\le f(x^{(n_k)})<m$ следует, что $f(x^{(n_k)})\to m$. Следовательно, $f(x^0)=m$, то есть максимум достигается. \\

        Аналогично для $m<0$ получим, что существует точка, в которой достигается минимум.
    \end{enumerate}
    
    Таким образом, $f$ ограничена на $K$ и достигает своих наибольшего и наименьшего значений.
\end{proof}

\definition{
$\{x^{(n)}\}$ \textit{равномерно сходится к точке} $x^0$, если $\forall\varepsilon>0\ \exists\delta >0\colon \forall n, n'\colon \rho(x^{(n)},x^{(n')})<\delta$, выполняется
\[
\rho(f(x^{(n)}),f(x^{(n')}))<\varepsilon.
\]
}

\theorem[Кантора]{Пусть $f\in C(K)$, $K$ --- компакт. \\
Тогда $f$ равномерно непрерывна на $K$.}
\begin{proof}
    От противного. Пусть равномерная непрерывность отсутствует. Тогда
    \[
    \exists\varepsilon_0>0\colon \forall\delta>0\ \exists x,y\in K\colon \rho(x,y)<\delta,\ \text{ выполняется } |f(x)-f(y)|\ge \varepsilon_0.
    \]
    Возьмём $\delta_n=\frac{1}{n}$. Тогда $\forall n\exists x^{(n)},y^{(n)}\in K \colon \rho(x^{(n)},y^{(n)})<\frac{1}{n}$, выполняется
    \[ 
    |f(x^{(n)})-f(y^{(n)})|\ge \varepsilon_0.
    \]
    
    Так как $K$ компактно, из последовательности $\{x^{(n)}\}$ можно выделить сходящуюся подпоследовательность $\{x^{(n_k)}\}$: $\lim\limits_{k\to\infty} x^{(n_k)}=x^0\in K$. \\
    Из последовательности $\{y^{(n_k)}\}$ выделим сходящуюся подпоследовательность $\{y^{(n_{k_l})}\}$: $\lim\limits_{l\to\infty} y^{(n_{k_l})}=y^0\in K$. \\
    
    Зафиксируем $\varepsilon>0$. В силу сходимости найдётся номер $L\colon \forall l\ge L$:
    \[
    \rho(x^{(n_{k_l})},x^0)<\frac{\varepsilon}{3},\qquad
    \rho(y^{(n_{k_l})},y^0)<\frac{\varepsilon}{3},\qquad
    \frac{1}{n_{k_l}}<\frac{\varepsilon}{3}.
    \]
    Тогда
    \[
    \rho(x^0,y^0)\le \rho(x^0,x^{(n_{k_l})})+\rho(x^{(n_{k_l})},y^{(n_{k_l})})+\rho(y^{(n_{k_l})},y^0)<
    \frac{\varepsilon}{3}+\frac{1}{n_{k_l}}+\frac{\varepsilon}{3}<\varepsilon.
    \]
    
    Так как $\varepsilon>0$ произвольно, то $x^0=y^0$. В силу непрерывности $f$ получим
    \[
    \lim_{l\to\infty} f(x^{(n_{k_l})})=\lim_{l\to\infty} f(y^{(n_{k_l})})=f(x^0).
    \]
    Следовательно, $\lim_{l\to\infty} |f(x^{(n_{k_l})})-f(y^{(n_{k_l})})|=0$. Значит, $\exists L^*\colon \forall l\ge L^*$ выполнено 
    \[ 
    |f(x^{(n_{k_l})})-f(y^{(n_{k_l})})|<\varepsilon_0.
    \]
    По построению $\forall l\colon |f(x^{(n_{k_l})})-f(y^{(n_{k_l})})|\ge \varepsilon_0$. Противоречие.
\end{proof}

\newlection{9 сентября 2026 г.}

\definition{
Пусть $f$ определена в точке $x^0\in E$. \\
Тогда $f$ \textit{непрерывна в точке} $x^0$, если
\[
\lim_{x\to x^0}f(x)=f(x^0).
\]
}

\definition{
Пусть $E\subset\mathbb R^d$, $f\colon E\to\mathbb R$, и $x^0$ --- предельная точка множества $E$. \\
Будем говорить, что $A=\lim\limits_{x\to x^0}f(x)$, если
\[
\forall \varepsilon>0\ \exists \delta>0\colon \forall x\in E\cap\mathring{U}_\delta(x^0),\ \text{выполняется}\ |f(x)-A|<\varepsilon,
\]
}

\theorem[5]{
Пусть $f$ определена на $E\subset \mathbb R^d$, $A\in \mathbb R$, $x^0$ --- предельная точка множества $E$. \\
Для того чтобы $\lim_{x\to x^0}f(x)=A$, необходимо и достаточно, чтобы для любой последовательности $\{x^{(n)}\}\colon x^{(n)}\to x^0$ выполнялось $f(x^{(n)})\to A$.
}
\begin{proof}
    \begin{enumerate}
        \item[$\Longrightarrow$] \textit{Необходимость.} Пусть $\lim_{x\to x^0}f(x)=A$. Тогда
        \[
        \forall \varepsilon>0\ \exists \delta>0\colon \forall x\in \mathring{U}_\delta(x^0),\ |f(x)-A|<\varepsilon.
        \]
        Возьмём последовательность $\{x^{(n)}\}\colon x^{(n)}\to x^0$. Тогда для данного $\delta$ найдётся номер $N\colon \forall n\ge N$ выполнено $\rho(x^{(n)},x^0)<\delta$, то есть $x^{(n)}\in \mathring{U}_\delta(x^0)$. \\
        Следовательно, $|f(x^{(n)})-A|<\varepsilon$.

       \item[$\Longleftarrow$] \textit{Достаточность}. Предположим, что для любой последовательности $\{x^{(n)}\}\colon x^{(n)}\to x^0$ выполняется $f(x^{(n)})\to A$. Докажем, что $\lim_{x\to x^0}f(x)=A$. \\
        Допустим противное:
        \[
        \exists \varepsilon_0>0\colon \forall \delta>0\ \exists x\in \mathring{U}_\delta(x^0),\ |f(x)-A|\ge \varepsilon_0.
        \]
        Положим $\delta_n=\frac1n$. Обозначим $y^{(n)}:=x(\delta_n)$. Тогда $\rho(y^{(n)},x^0)<\frac1n\to0$, значит, $y^{(n)}\to x^0$. \\
        По построению $|f(y^{(n)})-A|\ge \varepsilon_0$, что противоречит предположению о том, что для любой последовательности, сходящейся к $x^0$, значения функции сходятся к $A$.
    \end{enumerate}
\end{proof}



\end{document}

\end{document}


